\section{语言模型真的在推理吗？}

前面展示的各种方法在 benchmark 上取得了不错的成绩，但一个更深层的问题是：\textbf{语言模型是真正在推理，还是在利用记忆和模式匹配来模拟推理的表象？}

\begin{figure}[H]
\centering
\includegraphics[width=0.95\textwidth]{slides-images/slide-033.jpg}
\caption{``语言模型能推理吗？''——媒体报道充满矛盾：有人声称 AI 展现出人类推理的迹象，也有人认为 LLM 存在根本性的推理缺陷\protect\footnotemark}
\end{figure}
\footnotetext{Slides 第34页（标注为31）。}

\subsection{CoT 推理链的忠实性}

第一个问题：模型生成的推理链是否\textbf{忠实地}反映了其决策过程？

实验设计如下：让模型正常生成完整的推理链（比如4句话），然后在不同位置\textbf{强制截断}，让模型直接给出答案。如果推理链是忠实的，截断后答案应该会改变。

\begin{figure}[H]
\centering
\includegraphics[width=0.95\textwidth]{slides-images/slide-031.jpg}
\caption{Chain-of-Thought 忠实性实验：x轴为推理步骤的截断比例，y轴为准确率。在部分数据集上，即使大幅截断推理链，准确率几乎不变\protect\footnotemark}
\end{figure}
\footnotetext{Slides 第32页。}

\begin{warningbox}{推理链可能是``事后解释''}
实验发现：在多个数据集上，即使\textbf{完全截断推理链}（不让模型看到任何推理步骤），模型依然给出相同的答案。这意味着在这些情况下，推理链可能只是模型已确定答案后的\textbf{事后合理化解释}（post-hoc rationalization），而非真正驱动决策的推理过程。
\end{warningbox}

另一个实验进一步验证了这一点：在推理链的中间\textbf{故意插入错误}，观察是否影响最终答案。如果模型忠实地依赖推理链，插入错误应该导致答案改变。但实验发现，在部分数据集上，即使在第一步就引入错误，最终答案依然不变。

\subsection{反事实评估：推理还是记忆？}

第二个更尖锐的问题：模型在标准 benchmark 上的好表现，是真的因为掌握了推理能力，还是因为训练数据中包含了\textbf{类似的例子}？

\begin{figure}[H]
\centering
\includegraphics[width=0.95\textwidth]{slides-images/slide-037.jpg}
\caption{反事实实验——算术与逻辑：在标准设置（如十进制加法、符合常识的逻辑推理）下表现良好的模型，在反事实设置（如九进制加法、违反常识的逻辑前提）下性能显著下降\protect\footnotemark}
\end{figure}
\footnotetext{Slides 第38页（标注为34）。Source: Wu et al., 2024。}

实验设计了两种反事实：

\begin{enumerate}
    \item \textbf{算术反事实}：模型能做十进制加法（如 $12 + 14$），但十进制加法在训练数据中大量存在。如果改为\textbf{九进制加法}，模型的准确率\textbf{大幅下降}——说明它可能是在匹配模式而非真正理解进位规则。
    \item \textbf{逻辑反事实}：模型能正确处理``所有狗是哺乳动物''这类符合常识的逻辑推理。但如果改为\textbf{违反常识}的前提（如``柯基是爬行动物''），模型的推理能力显著退化——说明它可能依赖常识记忆而非纯逻辑推导。
\end{enumerate}

\subsection{类比推理中的反事实}

类比推理任务进一步揭示了这一问题。标准的类比推理（如扩展字母序列 abcd $\rightarrow$ abcde）语言模型可以完成得很好。

\begin{figure}[H]
\centering
\includegraphics[width=0.95\textwidth]{slides-images/slide-035.jpg}
\caption{类比推理的反事实测试：上半部分为标准字母表上的变换任务；下半部分使用\textbf{合成字母表}（xylkwbfz...）的变换任务，模型性能显著下降\protect\footnotemark}
\end{figure}
\footnotetext{Slides 第36页。Source: Hodel et al. 2024。}

但当引入两种反事实修改时：

\begin{itemize}
    \item \textbf{修改任务定义}：将``扩展序列''从输出下一个字符改为输出下下个字符
    \item \textbf{修改字母表}：将标准字母表替换为合成字母表（如 xylkwbfztn...）
\end{itemize}

模型性能\textbf{显著下降}。而在人类受试者上进行同样的实验，人类的准确率几乎不受影响。

\begin{importantbox}{推理 vs 记忆：目前的结论}
综合这些实验，目前的证据表明：
\begin{enumerate}
    \item 语言模型\textbf{确实展现了某种程度的推理能力}——它们能够在标准设置下解决多步推理问题
    \item 但这种推理能力\textbf{并不系统}——面对反事实或分布外的变体时，性能急剧下降
    \item 相当一部分``推理''可能是\textbf{模式匹配和记忆}的结果，而非真正的逻辑推导
\end{enumerate}
这一领域仍在快速发展，未来的研究可能会改变这些结论。
\end{importantbox}

\subsection{本章小结}

通过忠实性测试和反事实评估，我们发现语言模型的``推理''能力存在显著局限：推理链可能是事后合理化而非真正的推理过程，模型在偏离训练分布的问题上性能显著下降。这些发现提醒我们，在使用 LLM 进行推理时需要保持审慎态度。

%% ========== Part 5: 智能体基础 ==========

